% Verbatim copy of /Users/alingalatan/workspaces/grothe/handoff/manuscript/legacy/ch35-gpt-gaussian-remainder.tex (file dated 2026-10-25). Only this line was added.
% Consolidated source module (formerly Chapter 35). Owner: lower-bound/08-shared-structure.tex
\section{Lower-bound GPT: shared $A/B$ for arbitrary Gaussian-factor drivers, and a $15/16$ Gaussian Young remainder (20 October)}\label{ch:gaussian-remainder}
\chapsummary{Gaussian-factor drivers with coefficients independent of the shared Gaussian vector, plus the Gaussian Young remainder and covariance reduction. Independence of those coefficients is essential; nonlinear self-interaction is not covered by this theorem.}

\paragraph{Provenance and status (IDE note, 20 October).} Lower-bound GPT's \code{report/chapter\_new\_gaussian\_remainder.tex} and \code{content/gaussian\_marginal\_obstruction.tex}
from \pathref{gpt-reviews/k-next-20261020/}, ported verbatim apart from this header. \emph{Our replay} \st{C}: \code{python3 -I verify\_all.py} (standard library and Apple clang 17) from an
unrelated directory: \textbf{PASS full listed exact replay in 274\,s}. It covers the inverse-radius and $15/16$ / $43/45$ Young checks, the scope and tails, the vector Gram identity, the signed
obstructions, data regenerated byte for byte, the independent benchmark $S\ge96683367193980/2^{46}$, and the orientation cover: 11{,}799 splits, 10{,}400 accepted boxes, 1{,}400 tails, worst
$96686878076733$ against target $96686885631188$ (over $2^{46}$). Log \pathref{gpt-reviews/k-next-20261020/replay\_ide.log}.

% Proposed chapter fragment.  This file does not modify the live manuscript.
% Requires amsmath, amssymb, and hyperref; no custom theorem environments.
% Both complete arbitrary-driver covers have passed; the strongest retains the A/B sector penalty.
\subsection{A certified centered shared-sector theorem and a smaller Gaussian remainder}
\label{sec:new-gaussian-remainder}

\paragraph{Status and scope.}
The full arbitrary-driver centered $A/B$ cover has passed in exact
arithmetic, including the stronger sector-penalty target.  At the peak
profile this proves $K_\eta$ with $\eta=1/20000$ and a retained penalty
$(\|\alpha\|^2+\|\beta\|^2)/200$ for every centered Gaussian-factor
residual sharing any number of Gaussian coupling drivers, with no
amplitude restriction.
The proof uses a global $15/16$ Gaussian Young remainder, an exact Gram
reduction to three faces, and a convex orientation envelope that permits a
three-variable vertex certificate.  The scope is stated precisely in
Theorem~2.  General cubic self-interaction remains outside that theorem,
and the checked project lower bound remains
$K_G^{\mathbb C}>1.401674103\stale{22 Oct}{now $\KG>1.401695322458$}$.
In particular, this class satisfies full Gaussian domination whenever
its combined $A/B$ variance is at least $1/100$.

Throughout, a standard proper complex Gaussian has density
$\pi^{-1}e^{-|z|^2}$ and second moment $1$.  All Gaussian variables declared
independent are mutually independent.

\subsubsection{A global Gaussian norm inequality}

\paragraph{Theorem 1 (Gaussian Young remainder).}
Let $X,Y$ be independent standard proper complex Gaussians, and put
\[
 z=\ell+cX+gY,\qquad d=\beta\overline X,
 \qquad \ell,g\ge0,\quad c,\beta\in\mathbb C.
\]
If $z$ is not identically zero, then
\begin{equation}
 \boxed{\quad
 2\mathbb E|z+d|\le 2\mathbb E|z|
  +2\operatorname{Re}\mathbb E\frac{d\overline z}{|z|}
  +\frac{15}{16}|\beta|^2\mathbb E\frac{|X|^2}{|z|}.
 \quad}
 \label{eq:new-young}
\end{equation}
The values assigned to quotients on null sets are immaterial.  In
particular, no small-amplitude hypothesis on $\beta$ is present.

Here is the proof, including the finite polynomial certificate that supplies
the constant.  Scaling the variance of $z$ to $1$ and rotating $X,Y$ reduces
the positive-variance case to
\[
 z=r+aX+\sqrt{1-a^2}\,Y,\qquad d=u\lambda\overline X,
 \qquad r,u\ge0,\quad 0\le a\le1,\quad |\lambda|=1.
\]
The case $u=0$ is immediate.  Write
\[
 h=1+r^2,\quad F=h+a^2+u^2,\quad
 B=|z+d|^2-|z|^2,\quad C=\mathbb E\frac{|X|^2}{|z|}.
\]
The pointwise square discarded by Young's inequality gives the exact identity
\begin{align}
 2\mathbb E|z+d|
 &=2\mathbb E|z|+2\operatorname{Re}\mathbb E\frac{d\overline z}{|z|}
       +u^2C-\mathcal G,\\
 \mathcal G&=\mathbb E\frac{(|z+d|-|z|)^2}{|z|}.
 \label{eq:young-square}
\end{align}
With $D=|z|(|z+d|+|z|)^2$, weighted Cauchy--Schwarz yields
\begin{equation}
 \mathcal G=\mathbb E\frac{B^2}{D}
       \ge\frac{(\mathbb EB^2)^2}{\mathbb E(B^2D)}.
 \label{eq:weighted-square}
\end{equation}
The weights here are essential: the denominator is $\mathbb E(B^2D)$,
not $(\mathbb EB^2)(\mathbb ED)$.

Define the four polynomial moments
\[
 M_i=\mathbb E[B^2|z|^{2i}]\ (i=1,2),\qquad
 N_i=\mathbb E[B^2|z|^{2i}|d|^2]\ (i=0,1),
\]
and set $\tau=9h/4$ and
\[
 J=\frac{M_2+\tau M_1}{2\sqrt\tau},\qquad
 K=\frac{N_1+\tau N_0}{2\sqrt\tau}.
\]
The inequalities
\[
 2\sqrt\tau\,|z|^3\le |z|^4+\tau|z|^2,
 \qquad 2\sqrt\tau\,|z|\le |z|^2+\tau
\]
imply $\mathbb E[B^2|z|^3]\le J$ and
$\mathbb E[B^2|z||d|^2]\le K$.  Since
$D\le |z|(2|z|+|d|)^2$, Cauchy--Schwarz for its cross term gives
\begin{equation}
 \mathbb E(B^2D)\le (2\sqrt J+\sqrt K)^2.
 \label{eq:moment-denominator}
\end{equation}
Keeping this square is useful: replacing its cross term by a constant
linear combination of $J$ and $K$ loses part of the final improvement.
All of these moments are finite.  The inverse-radius term is integrable in
the nondegenerate complex Gaussian case, including when $Y$ is absent.

\subsubsection{The explicit moment and inverse-radius bounds}

For every phase $\lambda$ one has
\begin{equation}
 \mathbb EB^2=2u^2F.
 \label{eq:five-moments-first}
\end{equation}
For the remaining four moments their values at $\lambda=1$ are upper bounds.
This phase assertion does not assume the maximizing phase.  Before
substituting $g=\sqrt{1-a^2}$, expand in
$r,a,g,u,X,\overline X,Y,\overline Y$ and Laurent powers of $\lambda$.
Every coefficient in each of
$B^2|z|^2,B^2|z|^4,B^2|d|^2,B^2|z|^2|d|^2$ is nonnegative.
Gaussian integration deletes unmatched powers and multiplies the others by
nonnegative factorials.  The real value for $|\lambda|=1$ is therefore at
most the sum of these coefficients, which is the value at $\lambda=1$.
For $B^2$ itself, all surviving Laurent powers are zero, giving
\eqref{eq:five-moments-first} exactly.

Write the four upper bounds as
$M_i\le2u^2A_i$ and $N_i\le2u^4B_i$.  Direct expansion gives
\begin{align*}
 A_1={}&4a^2r^2+2a^2u^2+4a^2+4ar^2u
         +r^4+r^2u^2+4r^2+u^2+2,\\
 A_2={}&2a^4u^2+12a^3r^2u+11a^2r^4+8a^2r^2u^2
          +36a^2r^2+8a^2u^2+18a^2\\
       &+8ar^4u+24ar^2u+r^6+r^4u^2+9r^4+4r^2u^2
          +18r^2+2u^2+6,\\
 B_0={}&4a^2+2r^2+3u^2+2,\\
 B_1={}&4a^4+16a^2r^2+9a^2u^2+16a^2+12ar^2u
          +2r^4+3r^2u^2+8r^2+3u^2+4.
\end{align*}
Set
\[
 \mathsf A=A_2+\tfrac94hA_1,\qquad
 \mathsf B=B_1+\tfrac94hB_0.
\]
Then $J\le (2/3)u^2\mathsf A/\sqrt h$ and
$K\le (2/3)u^4\mathsf B/\sqrt h$.
Equations \eqref{eq:weighted-square}--\eqref{eq:five-moments-first} imply
\begin{equation}
 \mathcal G\ge
 \frac{6u^2\sqrt h\,F^2}{(2\sqrt{\mathsf A}+u\sqrt{\mathsf B})^2}.
 \label{eq:gap-rational}
\end{equation}

The second input is the pair of uniform inequalities
\begin{equation}
 H_0(r):=\mathbb E|r+X|^{-1}\le\frac{1773}{1000\sqrt{1+r^2}},
 \qquad
 H_1(r):=\mathbb E\frac{|X|^2}{|r+X|}
       \le\frac{38}{25\sqrt{1+r^2}}.
 \label{eq:normalized-inverse-radius}
\end{equation}
Gaussian regression then gives
\begin{equation}
 C=(1-a^2)H_0(r)+a^2H_1(r)
   \le\frac{C_*}{\sqrt h},\qquad
 C_*:=\frac{1773-253a^2}{1000}.
 \label{eq:C-regression}
\end{equation}

For completeness, these analytic functions and their certified domains are
explicit.  With $x=r^2/2$,
\begin{align*}
 H_0(r)&=\sqrt\pi e^{-x}I_0(x),\\
 H_1(r)&=\sqrt\pi e^{-x}
       \bigl[(x+\tfrac12)I_0(x)-xI_1(x)\bigr].
\end{align*}
For $H_0$, the ratio $R=I_1/I_0$ satisfies
$R'=1-R/x-R^2$.  The barrier $2x/(2x+1)$ starts strictly above $R$;
at a crossing the derivative of $R$ minus the barrier derivative is
$-3/(2x+1)^2$.  Hence $R<2x/(2x+1)$ for $x>0$ and
$\sqrt{1+2x}e^{-x}I_0(x)$ is decreasing from $1$.
The rational inequality $\sqrt\pi<1773/1000$ proves the first bound.

For $H_1$, positive power series for $I_0,I_1$, rational bounds for $\pi$,
and outward rational arithmetic verify the squared second bound on
$0\le x\le32$.  There are $32768$ cells of length $1/1024$.
The Bessel series are retained through index $128$, with their remaining
positive tail bounded geometrically.  This is an interval cover, not a
sample of point values.  On the unbounded tail $r\ge8$, the Laplace
integral substitution $t=s/(1+s)$ and a split at $t=1/2$ give
\begin{align}
 rH_0(r)&\le 1+\frac{1}{2r^2}
                       +\sqrt\pi r e^{-r^2/2},\\
 rH_1(r)&\le 1+\frac{3}{4r^2}+\frac{15}{8r^4}
             +\frac{\sqrt\pi}{2}r(1+r^2)e^{-r^2/2}.
 \label{eq:inverse-tails}
\end{align}
The polynomial and exponential terms decrease there.  Their endpoint
bounds prove \eqref{eq:normalized-inverse-radius} on the entire tail;
the replay uses $e>8/3$ and $\sqrt\pi<2$ for these comparisons.

\subsubsection{The finite polynomial certificate and completion of the theorem}

The remaining assertions are the two polynomial inequalities
\begin{align}
 P&:=96hF^2-C_*(4\mathsf A+u^2\mathsf B)\ge0,\\
 Q&:=P^2-16C_*^2u^2\mathsf A\mathsf B\ge0
 \qquad(r,u\ge0,\ 0\le a\le1).
 \label{eq:polynomial-gap-certificate}
\end{align}
Since $C_*,\mathsf A,\mathsf B$ are positive, these give
\[
 P\ge4C_*u\sqrt{\mathsf A\mathsf B},\qquad
 C_*(2\sqrt{\mathsf A}+u\sqrt{\mathsf B})^2\le96hF^2.
\]
Their global proof uses rational Bernstein coefficients after an explicit
compactification.  Both polynomials involve only even powers of $r$.
For either polynomial write $p(t,a,u)$ with $t=r^2$, and let its three
coordinate degrees be $(d_t,d_a,d_u)$.  Substitute
\[
 t=\frac{T_0}{1-T_0},\qquad u=\frac{U_0}{1-U_0},
\]
and multiply by $(1-T_0)^{d_t}(1-U_0)^{d_u}$.
For $T_0,U_0<1$ this multiplier is positive; the resulting expression is
polynomial on the whole cube $[0,1]^3$.
If $p=\sum p_{ijk}t^ia^ju^k$, its tensor Bernstein coefficient at
$(i,j,k)$ is exactly
\[
 \frac{1}{\binom{d_t}{i}\binom{d_u}{k}}
 \sum_{j'\le j}p_{ij'k}\frac{\binom j{j'}}{\binom{d_a}{j'}}.
\]
This formula follows directly from
$t^i(1-T_0)^{d_t}=T_0^i(1-T_0)^{d_t-i}$ after substitution, and the usual
power-to-Bernstein conversion in $a$.

The exact checker recursively bisects this cube, using rational de
Casteljau subdivision.  Both $P$ and $Q$ have nonnegative Bernstein
coefficients on every final box: there are $20$ boxes and maximum
subdivision depth $9$.  Hence they are nonnegative on the whole cube.
Compactification covers every finite $r,u\ge0$; it is not a truncation of
those unbounded parameters.

The replay reconstructs $P,Q$ from formal Gaussian moments using only
\[
 \mathbb E[X^j\overline X^{\,k}]
 =\begin{cases}j!,&j=k,\\0,&j\ne k,\end{cases}
\]
and verifies their coefficients before testing positivity.  It therefore
checks polynomial identities in the parameters and exact positivity
certificates, not sampled parameter values.

Finally, \eqref{eq:gap-rational}, \eqref{eq:C-regression}, and
\eqref{eq:polynomial-gap-certificate} give
\[
 \mathcal G\ge\frac{6}{96}\frac{u^2C_*}{\sqrt h}
          \ge\frac1{16}u^2C.
\]
Substitution into \eqref{eq:young-square} proves Theorem~1.  A deterministic
nonzero baseline follows by letting its Gaussian variance tend to zero.
The expected inverse-radius term has the corresponding continuous limit;
one can verify this by splitting outside a fixed neighborhood of the
nonzero deterministic shift and estimating the Gaussian density inside.

\paragraph{Dimension-free and conditional forms.}
The same conclusion holds for a standard proper Gaussian vector
$X\in\mathbb C^n$, with
\[
 z=\ell+\sum_jc_jX_j+gY,\qquad d=\sum_j\beta_j\overline X_j,
\]
provided the last term in \eqref{eq:new-young} is written as
$(15/16)\mathbb E[|d|^2/|z|]$.  A unitary change of coordinates puts
$d=\|\beta\|\overline X_1$; the remaining holomorphic Gaussian combination
is independent of $X_1$ and can be absorbed into $gY$.
Equivalently, this applies to any centered complex Gaussian pair $(Z,D)$
for which $(Z,\overline D)$ is jointly proper, with $z=\ell+Z,d=D$.
In elementary covariance terms the conditions are
$\mathbb EZ^2=\mathbb ED^2=\mathbb E[Z\overline D]=0$; the pairing
$\mathbb E[ZD]$ need not vanish.  The pair $(Z,D)$ itself need not be jointly
proper.  The result also applies conditionally on arbitrary latent data
when these hypotheses hold fibrewise and the displayed quantities are
integrable.  The linear pairing term must be retained under that
conditioning; independence of the latent variables does not make it zero.

\paragraph{Why the covariance condition matters.}
The improvement cannot hold for all Gaussian pairs.  Take $z=X$ and
$d=itX$, where $t$ is real.  The linear term vanishes, while the ratio of
the remaining norm increment to $\mathbb E[|d|^2/|z|]$ is
\[
 \frac{2(\sqrt{1+t^2}-1)}{t^2}
 =\frac{2}{\sqrt{1+t^2}+1}\longrightarrow1.
\]
Thus no universal coefficient below $1$ can replace Young's remainder
without an appropriate coupling constraint.  In this example
$\mathbb E[z\overline d]=-it\ne0$.  The antiholomorphic $B$ structure
satisfies the covariance condition; a coherent holomorphic $A$ perturbation
does not.  This is a structural distinction, not an amplitude restriction.

\subsubsection{A simultaneous shared A/B problem with three real variables}

We now apply the theorem at the peak profile
\[
 \epsilon=\frac{1891}{2500},\quad
 a_0=\frac{699}{1250},\quad b_0=\frac{179}{1250},\quad
 \ell(T)=\sqrt T(a_0+b_0T),\quad
 A=T-1,\quad B=T/\sqrt2,
\]
where $T$ is exponential with mean $1$.  Let
$h(w,\ell)=\mathbb E|\ell+\sqrt wX|$ and
\[
 S=\sup_{v\ge0}\left\{2\mathbb E_T h(v,\ell(T))-v/\epsilon\right\}.
\]
For deterministic $s,\alpha,\beta\in\mathbb C$ and $g\ge0$, consider
\begin{align}
 \mathcal D(s,\alpha,\beta,g)
 ={}&2\mathbb E|\ell+(s+\alpha A)X+\beta B\overline X+gY|\\
 &-\frac{|s|^2+g^2+|\alpha|^2+|\beta|^2}{\epsilon}
 +\frac75\operatorname{Re}(s\overline\alpha)
 +\frac57\operatorname{Re}(\beta s).
 \label{eq:centered-dual}
\end{align}
A uniform bound $\mathcal D\le S+\eta$ would apply by conditioning to
every centered latent $S_0$ independent of $X,Y,T$, with $s=S_0$ and
fixed $\alpha,\beta,g$.  The two support terms have expectation zero.
More generally the same statement holds whenever their two expectations
vanish.  Centering $S_0$ alone does not justify adaptive sector coefficients
$\alpha(S_0),\beta(S_0)$ without these extra cancellation conditions.

Define the scaled inverse-radius functions
\[
 H_0(w,\ell)=\mathbb E|\ell+\sqrt wX|^{-1},\qquad
 H_1(w,\ell)=\mathbb E\frac{|X|^2}{|\ell+\sqrt wX|}.
\]
The same exact integral checker proves
\begin{equation}
 H_j(w,\ell)\le\frac4{3\ell}\quad(j=0,1),
 \qquad \mathbb E_T\frac{B^2}{\ell}<\frac{39}{50}.
 \label{eq:B-price-integrals}
\end{equation}
The compact Bessel cover and the tails \eqref{eq:inverse-tails} prove the
first assertion.  Outward rational sums in $\sqrt T$, together with an
infinite Gaussian-tail bound, prove the second.

Rotate $X$ so that $\alpha=\sqrt r\ge0$, and write
\[
 v=|s|^2+g^2,\quad r=|\alpha|^2,\quad
 p=\operatorname{Re}(s\overline\alpha),\quad q=|s|^2.
\]
For $r>0$, feasibility says $p^2\le vr$ and $p^2/r\le q\le v$.
Let $w(T)=v+2pA+rA^2$.  The following quantities depend only on $v,r,p$:
\begin{align*}
 F_0&=2\mathbb E_T h(w,\ell)-(v+r)/\epsilon+\tfrac75p,\\
 k_w(T)&=-\frac{\sqrt\pi}{2\sqrt w}
           e^{-\ell^2/(2w)}I_1\bigl(\ell^2/(2w)\bigr),\\
 f&=\tfrac57+2\mathbb E_T[Bk_w],\qquad
 h_1=2\mathbb E_T[ABk_w],\\
 C(v,r,p,q)&=\mathbb E_T B^2\left[
 H_1(w,\ell)+\frac{v-q}{w}\{H_0(w,\ell)-H_1(w,\ell)\}\right].
\end{align*}
Use continuous zero-variance limits; isolated values of $T$ do not affect
the integrals.  Regression gives the formula for $C$ exactly.  Likewise
\[
 \mathbb E_{X,Y}\frac{\overline X\,\overline z}{|z|}
 =(s+\sqrt r A)k_w,
 \qquad z=\ell+(s+\sqrt r A)X+gY,
\]
which gives the linear coefficient
$\operatorname{Re}[\beta(sf+\sqrt r\,h_1)]$ in the dual.

Theorem~1, conditionally on $T$, and
\eqref{eq:B-price-integrals} show that its quadratic coefficient is at most
$(15/16)C$, where
\[
 0\le C<\frac{26}{25},\qquad
 \frac1\epsilon-\frac{15}{16}C>0.
\]
Maximizing the resulting quadratic upper bound over all complex $\beta$
therefore proves
\begin{equation}
 \boxed{\quad
 \mathcal D\le F_0+
 \frac{qf^2+rh_1^2+2pfh_1}
 {4\{\epsilon^{-1}-(15/16)C(v,r,p,q)\}}.
 \quad}
 \label{eq:three-variable-target}
\end{equation}
This eliminates the entire complex $B$ amplitude globally.

For fixed $v,r,p$, both numerator and denominator in
\eqref{eq:three-variable-target} are affine in $q$, with positive
denominator.  A linear-fractional function attains its maximum on an
interval at an endpoint.  It suffices to consider $q=v$ and $q=p^2/r$.
When $r=0$, use $q=v$ and $q=0$, the latter having zero gain.  Both branches
have the common coordinates
\[
 x\in\mathbb R,\quad y,\alpha\ge0,\qquad
 v=x^2+y^2,\quad r=\alpha^2,\quad p=x\alpha,
 \qquad w=(x+\alpha A)^2+y^2.
\]
The first branch has $s=x+iy,g=0$; the second has $s=x,g=y$.
Thus only two real three-variable inequalities remain.  This reduction
concerns the proved quadratic upper bound; it does not assert that the
original nonlinear payoff attains its maximum on these two faces.

\paragraph{Compactness and the certificate domain.}
The original dual needs no unbounded amplitude search.  With
$V=|s|^2+g^2+|\alpha|^2+|\beta|^2$, its two support terms have sum at most
$4V/5$, because
$\tfrac12\sqrt{(7/5)^2+(5/7)^2}<4/5$.
Also $\mathbb E\ell^2<19/25$.  Consequently
\[
 \mathcal D\le 2\sqrt{19/25+V}
              -(\epsilon^{-1}-4/5)V<27/20\qquad(V\ge10).
\]
The right side is decreasing on this tail.  This excludes $v+r\ge10$;
inside $v+r\le10$, the unrestricted quadratic maximization over $\beta$
in \eqref{eq:three-variable-target} is a safe upper bound.

The complete upper bound in \eqref{eq:three-variable-target} has not been
proved convex.  A correct box certificate may use convexity only for the
norm part of $F_0$.  In the coordinates $x,y,\alpha$, the support term
$(7/5)x\alpha$ is multilinear and interpolates exactly.  Thus a boxwise
upper bound for $F_0$ is its largest certified vertex value plus
\[
 \frac{(\Delta x)^2+(\Delta y)^2+(\Delta\alpha)^2}{4\epsilon}.
\]
One then adds a direct interval enclosure of each gain in
\eqref{eq:three-variable-target}.  Directly interpolating the whole upper bound would not be justified.
Section~\ref{sec:convex-orientation-envelope} gives a different, valid
vertex method: it interpolates a convex envelope of the original norm,
using this upper bound only at the vertices.

The preceding version with coefficient $1$ had a genuine secondary maximum
in its quadratic upper bound: numerical optimization gives approximately
$1.37658$, above $S\simeq1.37395353$.  This was a failure of that upper bound,
not a counterexample to K, and subdivision alone could not fix it.
The factor $15/16$ addresses precisely this discarded-square loss.  The
complete interval cover is supplied below, including the additional
vector face, and proves slack $1/20000$.  General cubic self-interaction
remains outside its hypotheses.

\subsubsection{Arbitrarily many Gaussian drivers: a Gram reduction to three faces}
\label{sec:vector-gram-reduction}

The preceding scalar reduction extends to any number of Gaussian coupling
drivers.  The extension retains the quadratic coefficient as a matrix;
replacing it by a scalar bound would discard the direction of the linear
content term.

Let $X\in\mathbb C^n$ be standard proper Gaussian, let $Y$ be independent,
and let $s,\alpha,\beta\in\mathbb C^n$ be deterministic.  Use the same
$A,B,\ell,\epsilon$ as above, set $c(T)=s+A\alpha$, and define
\begin{align*}
 \mathcal D_n={}&2\mathbb E|\ell+c(T)^TX+B\beta^T\overline X+gY|\\
 &-\frac{\|s\|^2+g^2+\|\alpha\|^2+\|\beta\|^2}{\epsilon}
   +\frac75\operatorname{Re}(s^*\alpha)
   +\frac57\operatorname{Re}(\beta^Ts).
\end{align*}
A uniform bound on this deterministic dual can again be averaged over
latent coefficient vectors.  A centered random $s$, independent of
$X,Y,T$, cancels both support terms when $\alpha,\beta$ are fixed.  For
adaptive coefficients one must instead require the expectations of both
displayed support pairings to vanish.

Write $\theta=15/16$, $w=\|c(T)\|^2+g^2$ and
\[
 v=\|s\|^2+g^2,\qquad r=\|\alpha\|^2,\qquad
 p=\operatorname{Re}(s^*\alpha).
\]
The baseline $F_0$ and the real numbers $f,h_1$ remain exactly those in
\eqref{eq:three-variable-target}: they depend only on $v,r,p$, since
$w=v+2pA+rA^2$.  For the optimization variable $b=\overline\beta$,
Gaussian regression gives the Hermitian quadratic coefficient
\begin{equation}
 M=\mathbb E_T B^2\left[
 H_0(w,\ell)I+
 \frac{H_1(w,\ell)-H_0(w,\ell)}w c(T)c(T)^*\right].
 \label{eq:vector-M}
\end{equation}
Indeed the underlying conditional second-moment formula is
\[
 \mathbb E\frac{\overline X X^T}{|\ell+c^TX+gY|}
 =H_0(w,\ell)I+
   \frac{H_1(w,\ell)-H_0(w,\ell)}w cc^*.
\]
It follows that $0\le M\le(26/25)I$.  In particular the price matrix
\[
 Q=\epsilon^{-1}I-\theta M
\]
is strictly positive definite.  The dimension-free version of Theorem~1
then yields
\[
 \mathcal D_n\le F_0+operatorname{Re}(b^*L)-b^*Qb,
 \qquad L=sf+\alpha h_1.
\]
Completing the square in $b$ proves the exact upper reduction
\begin{equation}
 \boxed{\qquad
 \mathcal D_n\le F_0+\tfrac14L^*Q^{-1}L.
 \qquad}
 \label{eq:vector-upper}
\end{equation}
There is no dimension-dependent estimate in this step.

To reduce its remaining matrix term, introduce
\[
 C_0=\mathbb E_T B^2H_0(w,\ell),\quad
 d_j=\mathbb E_T\left[B^2A^j
              \frac{H_1(w,\ell)-H_0(w,\ell)}w\right]\quad(j=0,1,2),
\]
\[
 K=\begin{pmatrix}d_0&d_1\\d_1&d_2\end{pmatrix},\qquad
 \lambda=\epsilon^{-1}-\theta C_0>0,\qquad V=[s\ \alpha].
\]
For nonsingular variance profiles these are ordinary integrals and
$M=C_0I+VKV^*$.  At singular profiles the formulas below are understood
through the original integral matrix \eqref{eq:vector-M}, or by a limit.
Individual $d_j$ bounds near a zero of $w$ are not needed for the cover.
The span of $s,\alpha$ is invariant under $Q$ and contains $L$.
Its Gram matrix is
\[
 \mathcal G=V^*V=\begin{pmatrix}q&\kappa\\\overline\kappa&r\end{pmatrix},
 \quad q=\|s\|^2,\quad p=\operatorname{Re}\kappa,
 \quad \Delta=qr-|\kappa|^2.
\]
The two-by-two adjugate identity
\[
 \operatorname{adj}(\lambda I-\theta\mathcal GK)\mathcal G
       =\lambda\mathcal G-\theta\Delta\operatorname{adj}K
\]
therefore gives
\begin{equation}
 L^*Q^{-1}L=
 \frac{\lambda(qf^2+rh_1^2+2pfh_1)
       -\theta\Delta(d_2f^2-2d_1fh_1+d_0h_1^2)}
 {\lambda^2-\lambda\theta(d_0q+d_2r+2d_1p)
                  +\theta^2(d_0d_2-d_1^2)\Delta}.
 \label{eq:Gram-fraction}
\end{equation}
For example, this follows by writing $L=V(f,h_1)^T$, applying the
resolvent identity on the range of $V$, and multiplying its two-by-two
inverse by $\mathcal G$.  The denominator is
$\det(\lambda I-\theta\mathcal GK)>0$.  One way to see its sign is to
replace $\mathcal GK$ by the similar Hermitian expression
$\mathcal G^{1/2}K\mathcal G^{1/2}$ on the positive-rank subspace; its
price is the restriction of $Q>0$.  Rank-deficient cases follow directly
or by continuity, with the remaining eigenvalue equal to $\lambda>0$.

Fix $v,r,p$ with $r>0$.  All of $C_0,K,f,h_1,\lambda$ are now fixed.
The feasible set for the two remaining invariants is exactly the triangle
\[
 p^2/r\le q\le v,\qquad 0\le\Delta\le qr-p^2.
\]
Indeed $g^2=v-q$ and
$(\operatorname{Im}\kappa)^2=qr-p^2-\Delta$; any positive semidefinite
Gram matrix with these data is realized in complex dimension at most two.
The numerator and denominator of \eqref{eq:Gram-fraction} are affine on
this triangle, and the denominator is positive.  Such a ratio is bounded
by its vertex values: if a point has barycentric weights $t_i$, the ratio
there equals the weighted average of vertex ratios with positive weights
proportional to $t_i$ times the vertex denominators.  Hence only the three
vertices
\[
 (q,\Delta)=(p^2/r,0),\qquad (v,0),\qquad (v,vr-p^2)
\]
need be covered.  The first two are the previous scalar faces.  The third
is a new real two-driver face.  The case $r=0$ follows by a direct
one-vector calculation or continuity.

Use common coordinates $x\in\mathbb R$, $y,a\ge0$ with
$v=x^2+y^2$, $r=a^2$, $p=xa$.  The new face has
\[
 s=(x,y),\quad\alpha=(a,0),\quad g=0,\qquad
 c_1=x+aA,\quad c_2=y,\quad w=c_1^2+y^2,
 \quad L=(xf+ah_1,yf).
\]
For numerical certification the stable formula is the original matrix
integral:
\begin{align*}
 M_{11}&=\mathbb E_T B^2\left[H_0+(H_1-H_0)c_1^2/w\right],\\
 M_{22}&=\mathbb E_T B^2\left[H_0+(H_1-H_0)y^2/w\right],\\
 M_{12}&=\mathbb E_T B^2(H_1-H_0)c_1y/w.
\end{align*}
This avoids separate estimates of potentially singular $d_j$.
The same $v+r\ge10$ tail exclusion and the same convex vertex bound for
$F_0$ apply.  Only one additional three-variable matrix gain needs to be
certified, independently of the original driver dimension.

There is a useful enclosure rule for this gain.  If $U_0,U_1$ are
pointwise upper bounds for $H_0,H_1$, respectively, then
\[
 H_0P_\perp+H_1P_\parallel
 \le U_0P_\perp+U_1P_\parallel
\]
in positive-semidefinite order, where
$P_\parallel=cc^*/w$ on this face.  This holds regardless of the sign of
$U_1-U_0$.  Integrate an interval enclosure of the upper matrix.  Fixed
upper diagonal entries, together with its actual enclosed off-diagonal
entry, still dominate the true matrix.  Provided the resulting price
matrix stays positive throughout the interval, its inverse quadratic can
be bounded by the two off-diagonal endpoints and the four vertices of a
rectangular enclosure of the real vector $L$.  This follows from convexity
of $L^*Q^{-1}L$ in $Q>0$ for fixed $L$, and convexity in $L$ for fixed $Q$.
It does not require convexity of the complete reduced objective.

This proves the arbitrary-driver reduction.  Its numerical conclusion
comes from the complete exact cover reported below, including all three
faces, rather than exploratory optimization of \eqref{eq:vector-upper}.
Non-Gaussian cubic self-interaction remains outside the Gaussian-factor
hypotheses.

\subsubsection{A convex orientation envelope gives a genuine vertex certificate}
\label{sec:convex-orientation-envelope}

There is a better use of convexity than enclosing the reduced matrix gain
over an entire parameter box.  The original orientation envelope is convex
before subtracting the baseline price.  The quadratic majorant need only
supply certified values at the vertices; no convexity of that majorant is
required.

Let $X_1,X_2,X_3,Y$ be independent standard proper complex Gaussians.  For
$x\in\mathbb R$, $R,a\ge0$, and
$\omega=(\omega_1,\omega_2,\omega_3)\in\mathbb R_{\ge0}^3$ with
$\|\omega\|=1$, set
\[
 s_\omega=(x+iR\omega_1,R\omega_2,0),\qquad
 \alpha=(a,0,0),\qquad g=R\omega_3.
\]
Define the common envelope
\begin{align*}
 \mathcal H(x,R,a)=\sup_{\omega,\,\beta\in\mathbb C^3}
 \biggl\{&2\mathbb E\bigl|\ell+(x+iR\omega_1+aA)X_1
                +R\omega_2X_2+R\omega_3Y+B\beta^T\overline X\bigr|\\
 &-\|\beta\|^2/\epsilon
                    +\tfrac57\operatorname{Re}(\beta^Ts_\omega)\biggr\}.
\end{align*}
The supremum is finite: for fixed $x,R,a$ the norm and support terms grow
at most linearly in $\|\beta\|$, whereas its price is a negative quadratic.
For every fixed $\omega,\beta$, the expression inside the braces is convex
in the three real variables $(x,R,a)$.  It is an expected norm of an affine
expression, plus an affine support term and a constant price.  A supremum
of convex functions is convex, so $\mathcal H$ is convex on
$\mathbb R\times[0,\infty)^2$.

This envelope covers every original driver dimension.  A unitary change
makes $\alpha=(a,0,\ldots)$, and a further unitary change fixing that
coordinate makes $s=(x+iz,y,0,\ldots)$ with $y\ge0$.  Simultaneous complex
conjugation permits $z\ge0$ without changing the dual, and $g$ can be taken
nonnegative.  Put $R=\sqrt{z^2+y^2+g^2}$ and use its three normalized
components as $\omega$.  The part of $\beta$ orthogonal to the first two
coordinates is one independent proper Gaussian term; it can be combined
into the third coordinate with the same coefficient norm.  Its support
pairing with $s$ is zero.  Adding unused coordinates in dimensions below
three only enlarges the supremum, which is safe.  Consequently every
original dual configuration is bounded by
\[
 \widehat{\mathcal D}(x,R,a)
 =\mathcal H(x,R,a)-\frac{x^2+R^2+a^2}{\epsilon}+\frac75xa.
\]
The orientation-independent baseline price is essential:
$\|s_\omega\|^2+g^2+\|\alpha\|^2=x^2+R^2+a^2$, and
$\operatorname{Re}(s_\omega^*\alpha)=xa$.

At each fixed $(x,R,a)$, the Gram invariants are
\[
 v=x^2+R^2,\quad r=a^2,\quad p=xa,\qquad
 q=x^2+R^2(\omega_1^2+\omega_2^2),\quad
 \Delta=a^2R^2\omega_2^2.
\]
The squared components of $\omega$ are precisely barycentric coordinates
of the Gram triangle.  Its three vertices are
\[
 \omega=(0,0,1),\qquad(1,0,0),\qquad(0,1,0),
\]
representing respectively the scalar real driver plus buffer, scalar
complex driver, and real two-driver face.  Let $U_1,U_2,U_3$ be their
pointwise upper bounds from the Gaussian Young and Gram calculation.
The preceding reduction proves
\[
 \widehat{\mathcal D}(x,R,a)\le
 U(x,R,a):=\max\{U_1(x,R,a),U_2(x,R,a),U_3(x,R,a)\}.
\]
The three $U_i$ must bound the Gram-reduced quadratic face expressions,
not merely the actual norm payoffs at three orientations.  The latter
would not control interior orientations: no axis-extremality assertion is
made for the original nonlinear payoff.  This includes optimization over
every $\beta$ direction, including the third unused baseline coordinate.  No assertion that the original
nonlinear optimum has zero perpendicular $\beta$ is needed.

Now let $\mathcal B$ be a rectangular box with side lengths
$d_1,d_2,d_3$, and let $\lambda_v$ be the standard multilinear interpolation
weights of a point $z\in\mathcal B$ over its eight vertices $v$.  Convexity
of $\mathcal H$ gives
$\mathcal H(z)\le\sum_v\lambda_v\mathcal H(v)$.
The product $xa$ interpolates exactly.  For each coordinate $j$,
\[
 \sum_v\lambda_vv_j^2-z_j^2
 =t_j(1-t_j)d_j^2\le d_j^2/4,
\]
where $t_j\in[0,1]$ is the coordinate's relative location in the box.
It follows that
\begin{equation}
 \boxed{\quad
 \widehat{\mathcal D}(z)
 \le \max_{v\text{ a vertex of }\mathcal B}U(v)
          +\frac{d_1^2+d_2^2+d_3^2}{4\epsilon}.
 \quad}
 \label{eq:orientation-vertex-certificate}
\end{equation}
The values $U(v)$ may be any rigorously computed point upper bounds.
They need not be convex or even continuous as functions of the parameters.
The proof interpolates $\mathcal H$, corrects for the quadratic price, and
only then substitutes the majorants at vertices.  It never interpolates
$U$ itself.

For the scalar problem, the same argument uses a positive unit circle of
orientations and its two scalar faces.  For arbitrarily many drivers,
\eqref{eq:orientation-vertex-certificate} supplies a three-dimensional
vertex scheme with only one additional point evaluation per vertex.
The existing tail proof excludes $x^2+R^2+a^2\ge10$ for the original dual,
so $[-4,4]\times[0,4]^2$ covers the remaining domain.  On the fixed
$2^{-46}$ lattice, the outward curvature correction is exactly
\[
 \left\lceil
 \frac{625\sum_j d_{j,\mathrm{integer}}^2}{1891\,2^{46}}
 \right\rceil
\]
in integer output units.  Certified vertex values can therefore be cached
and reused across adjacent boxes; a full parameter-box enclosure of the
gain is unnecessary.  Radial integration and its unbounded tails must
still be certified at every evaluated vertex.

This proved computational reduction is the mechanism used in the
completed cover below.  The Gaussian-factor restriction remains: the
vertex argument does not by itself compare a general non-Gaussian cubic
self-interaction to the benchmark.

\subsubsection{Certified centered A/B theorem, uniform in the driver dimension}
\label{sec:certified-centered-AB}

The complete cover of all three Gram faces has passed, including the
stronger simultaneous sector penalty.  It proves the following form of K
for the Gaussian-factor class.

\paragraph{Theorem 2 (all centered Gaussian-factor drivers).}
Let $n\ge1$, let $X\in\mathbb C^n$ and $Y\in\mathbb C$ be independent
standard proper Gaussians, and let $T$ be an independent exponential
variable of mean $1$.  Let $S_0\in\mathbb C^n$ be any centered,
square-integrable random vector independent of $X,Y,T$.  For arbitrary
fixed $\alpha,\beta\in\mathbb C^n$ and $g\ge0$, put
\[
 R=(S_0+(T-1)\alpha)^TX+\frac{T}{\sqrt2}\beta^T\overline X+gY.
\]
At the profile and price fixed above,
\begin{equation}
 \boxed{\begin{aligned}
 &2\mathbb E|\ell(T)+R|
 -\frac{\mathbb E\|S_0\|^2+g^2+\|\alpha\|^2+\|\beta\|^2}{\epsilon}\\
 &\hspace{20mm}\le S+\frac1{20000}
            -\frac{\|\alpha\|^2+\|\beta\|^2}{200}.
 \end{aligned}}
 \label{eq:certified-centered-AB}
\end{equation}
There are no restrictions on the dimension, amplitudes, or law of $S_0$
beyond the hypotheses stated.  The same result holds for random
$S_0,\alpha,\beta,g$ independent of $X,Y,T$, with the necessary square
integrability and
\[
 \mathbb E\operatorname{Re}(S_0^*\alpha)=0,
 \qquad \mathbb E\operatorname{Re}(\beta^TS_0)=0.
\]
In that version the entire price and the sector penalty in
\eqref{eq:certified-centered-AB} are averaged.  The two cancellations, rather than centering alone, are the
condition needed for adaptive coefficients.

A useful immediate corollary is full Gaussian domination on this class
whenever $\|\alpha\|^2+\|\beta\|^2\ge1/100$: then the right-hand side
of \eqref{eq:certified-centered-AB} is at most $S$.  For adaptive random
coefficients, the same conclusion follows when their expected combined
sector variance is at least $1/100$.

The theorem also holds for countably many drivers, provided
$S_0$ is a centered $\ell^2$-valued random vector with
$\mathbb E\|S_0\|_{\ell^2}^2<\infty$ and
$\alpha,\beta\in\ell^2$.  Truncate all three vectors at coordinate $n$.
The residuals converge in $L^2$, by independence and the Gaussian isometry,
and their prices and sector penalties converge.  The inequality
$||z+r|-|z+r'||\le|r-r'|$ passes the two norm expectations to the limit.
This gives the same constants without an additional compactness assumption.

To prove the theorem, apply the deterministic dual bound with the
sector-price modification proved below and its completed orientation-envelope
certificate to each realization of $S_0$, and then average.  The two support pairings in the
dual vanish by centering, or by the stated cancellation conditions.
All integrals pass through this conditioning because the norm has linear
growth and the coefficients are square-integrable.

\paragraph{Original visible-coordinate notation.}
Equivalently, let $Z$ be an independent standard proper Gaussian and set
\[
 L(Z)=Z(a_0+b_0|Z|^2),\qquad
 \widetilde R=S_0^TX+gY+(|Z|^2-1)\alpha^TX
                       +\frac{Z^2}{\sqrt2}\beta^T\overline X.
\]
Then Theorem~2 states
\[
 \mathbb E|L+\widetilde R|+\mathbb E|L-\widetilde R|
 -\frac{\mathbb E|\widetilde R|^2}{\epsilon}
 \le S+\frac1{20000}-\frac{\|\alpha\|^2+\|\beta\|^2}{200}.
\]
Indeed write $Z=\sqrt T e^{i\vartheta}$, factor $e^{i\vartheta}$ out of
the norm, and rotate $X,Y$ by $e^{-i\vartheta}$.  Their joint conditional
Gaussian law is unchanged, and the two relative frequencies give precisely
the $T$-model above.  The conditional law of $R$ is symmetric, so the two
norm expectations coincide.  Finally
$\mathbb E|\widetilde R|^2=\mathbb E\|S_0\|^2+g^2+
\|\alpha\|^2+\|\beta\|^2$, by $\mathbb E(T-1)=0$,
$\mathbb E(T-1)^2=\mathbb ET^2/2=1$, and properness of the Gaussian drivers.

Here are the exact computational endpoints of the deterministic bound.
The independently certified Gaussian benchmark satisfies
\[
 S\ge S_{\rm lb}:=\frac{96683367193980}{2^{46}}.
\]
The full orientation cover used the target integer
\[
 C_{\rm target}=96686885631188
 =96683367193980+\left\lfloor\frac{2^{46}}{20000}\right\rfloor.
\]
Its final receipt is
\[
 \begin{array}{lr}
 \text{splits}&11799\\
 \text{accepted compact boxes}&10400\\
 \text{boxes discharged by the analytic tail}&1400\\
 \text{distinct certified vertices}&17505\\
 \text{largest accepted integer upper bound}&96686878076733.
 \end{array}
\]
Thus the largest certified value is strictly below the target.  The
compact boxes and tail boxes exhaust the initial cube; their number is
$11799+1$, as required by a completed binary subdivision.  The analytic
variance tail covers everything outside the compact region.  Equations
\eqref{eq:orientation-vertex-certificate} and the exact radial enclosures
therefore give
\[
 \mathcal D_n+\frac{\|\alpha\|^2+\|\beta\|^2}{200}
 \le \frac{C_{\rm target}}{2^{46}}
 \le S_{\rm lb}+\frac1{20000}\le S+\frac1{20000}
\]
uniformly in all deterministic coefficients.  This is the bound used in
the conditioning argument, not an interpolation of floating-point values.

This theorem closes simultaneous shared $A/B$ coupling for every centered
Gaussian-factor residual, with arbitrary dependence among its latent
coefficients.  It does not cover a general residual that is a nonlinear
cubic in the same Gaussian drivers, such as a coherent Wick self-interaction.
Consequently it does not prove full K or supply a new complete row
envelope.  The current checked project lower bound remains
$K_G^{\mathbb C}>1.401674103\stale{22 Oct}{now $\KG>1.401695322458$}$.

\subsubsection{The certified simultaneous sector penalty}
\label{sec:certified-sector-penalty}

The stronger certificate retains a simultaneous penalty for the two
shared sectors.  With $\gamma=1/200$, its target is
\begin{equation}
 \mathcal D_n+\gamma(\|\alpha\|^2+\|\beta\|^2)
 \le S+\frac1{20000}.
 \label{eq:bonus-target}
\end{equation}
This claim uses a separate exact cover, rather than an inference from
the unstrengthened cover.  The following changes give the analytic
reduction used by that stronger replay.

Write $p_0=\epsilon^{-1}$.  The quadratic $B$ optimization changes only by
\[
 Q_\gamma=(p_0-\gamma)I-\theta M,
 \qquad \lambda_\gamma=p_0-\gamma-\theta C_0,
 \qquad F_{0,\gamma}=F_0+\gamma a^2.
\]
The pointwise inequality
$p_0-\gamma-\theta(26/25)>0$ keeps this price positive.
The linear content vector does not change.  The numerator and denominator
of the Gram fraction remain affine in $(q,\Delta)$, so the same three-face
reduction is valid with $Q_\gamma$.

Define $\mathcal H_\gamma$ by changing the $B$ price in the orientation
envelope to $-(p_0-\gamma)\|\beta\|^2$.
Then $\mathcal H_\gamma+\gamma a^2$ is convex, and the strengthened original
envelope equals
\[
 (\mathcal H_\gamma+\gamma a^2)
       -p_0(x^2+R^2+a^2)+\tfrac75xa.
\]
Thus the same curvature correction
$p_0(d_1^2+d_2^2+d_3^2)/4$ remains valid.  One could use the slightly
smaller correction
$[p_0(d_1^2+d_2^2)+(p_0-\gamma)d_3^2]/4$ by interpolating only
$\mathcal H_\gamma$, but it is unnecessary for the validity of the
certificate.

For the unbounded tail, let $V$ be the full priced variance, including the
$B$ coefficient.  The strengthened dual is at most
\[
 2\sqrt{19/25+V}-(p_0-4/5-\gamma)V<27/20\qquad(V\ge11).
\]
This follows from the exact comparisons
\[
 \left(\frac{27}{20}+11\left(p_0-\frac45-\frac1{200}\right)\right)^2
       -4\left(\frac{19}{25}+11\right)
 =\frac{355784779281}{143035240000}>0,
\]
and $(p_0-4/5-1/200)^2(19/25+11)>1$, which proves that the upper bound is
decreasing on the tail.  The compact region changes from baseline variance
at most $10$ to at most $11$; the same initial cube still contains it.

The strengthened cover passed with the receipt in
Section~\ref{sec:certified-centered-AB}.  Conditioning removes the support
terms exactly as before and leaves the additional deficit
$-\gamma\mathbb E(\|\alpha\|^2+\|\beta\|^2)$.
No penalty for the remote latent residual variance or the independent
Gaussian buffer is asserted.  Theorem~2 includes precisely this certified
sector penalty.

\subsubsection{Why a general replacement must retain more information}

Two rigorous examples constrain the next step.  They do not refute K.

\paragraph{Visible averaging does not make the signed replacement favorable.}
Let $H(X)=X(|X|^2-2)/\sqrt2$ and
$R_T=cH(X)+\alpha(T-1)X$, with real $c,\alpha$.  Conditional on $T$, this is
circular with variance $v_T=c^2+\alpha^2(T-1)^2$.  Its fourth cumulant is
\[
 \kappa_T=42c^4+32\sqrt2c^3\alpha(T-1)
     +20c^2\alpha^2(T-1)^2+4\sqrt2c\alpha^3(T-1)^3.
\]
At $c=-\alpha/(4\sqrt2)$, its visible average equals
$-683\alpha^4/512$.  A conditional slice contains first and third remote
chaos, so this sign is compatible with a pure total cubic.
If $\Gamma_T$ is proper Gaussian of variance $v_T$, an exact Rice Taylor
remainder check gives, at $\alpha=1$ and buffer $g=32$,
\[
 2\mathbb E|\ell+R_T+32Y|
 -2\mathbb E|\ell+\Gamma_T+32Y|>\frac{19}{10^7}>0.
\]
This is an admissible shared-$A$ construction (the independent graded
Gaussian can be approached by finite remote Wick sums).  The matched
Gaussian profile has its own strict amplitude penalty and lies far from
the optimizing variance here.  Thus the correct comparison must keep that
penalty; a universal nonpositive signed error would be false.

\paragraph{An exactly Gaussian radial law still does not determine shared content.}
Let $v_*>0$ maximize the Gaussian benchmark and define
\[
 R=|X|^2,\qquad F_2(r)=1-(r+1)e^{-r},\qquad
 Q=\sqrt{v_*}\,X e^{2\pi iF_2(R)}.
\]
The variable $Q$ is exactly proper Gaussian of variance $v_*$: it preserves
the radius and translates a uniform phase.  Moreover,
\[
 \mathbb E[Q\overline X]/\sqrt{v_*}
 =\int_0^\infty re^{-r}e^{2\pi iF_2(r)}\,dr
 =\int_0^1e^{2\pi iy}\,dy=0.
\]
The conditional covariance profile of $Q+\alpha(T-1)X$ is therefore
$v_*+|\alpha|^2(T-1)^2$, with zero pseudovariance, exactly as in the
Gaussian $A$ comparison.

Nevertheless, some arbitrarily small imaginary $A$ perturbations raise
the payoff above $S$.  The previously derived averaged $A$ potential
$F_A(u)=\mathbb E[(T-1)|L+x|]$, $|x|^2=u$, has
$w(u)=-F_A'(u)>0$ strictly decreasing, by its positive Laplace mixture.
For
$J=\mathbb E[R e^{2\pi iF_2(R)}w(v_*R)]$, pairing $y$ with $1-y$ gives
\[
 \operatorname{Im}J=\int_0^{1/2}\sin(2\pi y)
  [w(v_*F_2^{-1}(y))-w(v_*F_2^{-1}(1-y))]\,dy>0.
\]
The first variation is
$-4\sqrt{v_*}\operatorname{Re}(\alpha\overline J)$, so
\[
 \Lambda(Q-it(T-1)X)=S+4\sqrt{v_*}(\operatorname{Im}J)t+o(t)>S
 \quad(t>0\text{ sufficiently small}).
\]
This example is not pure cubic.  In real standard Gaussian coordinates its
Dirichlet energy is
\[
 \mathbb E\|DQ\|^2=v_*\left(1+\frac{64\pi^2}{81}\right)>3v_*,
\]
whereas a pure cubic has energy exactly $3v_*$.  Exact Gaussian marginals,
charge one, first-chaos orthogonality, and the correct covariance profile
therefore do not replace the joint fixed-chaos identities.

\paragraph{A caution about fresh-noise tests.}
For a Gaussian-smoothed cubic objective, let
$e=\Pi_{2,1}k_g-G/\epsilon$ be its first-variation residual.  The quadratic
form for a fresh polarized noise $W_1P+\overline{W_1}Q$,
$P\in\mathcal H_{1,1},Q\in\mathcal H_{2,0}$, has polarization
$\mathfrak b_g$ satisfying, for each remote coordinate $j$,
\[
 \mathfrak b_g((D_jG,\overline D_jG),(P,Q))
 =\operatorname{Re}\mathbb E[
       (D_je)\overline P+(\overline D_je)\overline Q].
\]
Thus derivative probes form a forced nullspace at every stationary point.
Their trace only repeats stationarity.  A useful extra noise-stability
constraint must retain probes outside that nullspace, together with their
paired $(1,1)/(2,0)$ structure.

\subsubsection{Replay and handoff}

The principal standard-library exact checkers are located at the following
absolute paths (run each with \texttt{python3 -I}):
\begin{quote}\small\ttfamily\raggedright
\nolinkurl{/Users/alingalatan/Documents/Codex/2026-10-10/plea/outputs/k-next-20261020/dual/check_inverse_radius.py}\\[3pt]
\nolinkurl{/Users/alingalatan/Documents/Codex/2026-10-10/plea/outputs/k-next-20261020/sharpen/check_gaussian_young_15_16.py}\\[3pt]
\nolinkurl{/Users/alingalatan/Documents/Codex/2026-10-10/plea/outputs/k-next-20261020/gamma/check_signed_obstructions.py}
\end{quote}
The Gaussian-gap replay rebuilds five formal Gaussian moments and the two
polynomials, and proves their positivity on $20$ compactified Bernstein
boxes.  The inverse-radius replay supplies its separate analytic-function
inputs.  The stronger orientation cover is recorded at
\begin{quote}\small\ttfamily\raggedright
\nolinkurl{/Users/alingalatan/Documents/Codex/2026-10-10/plea/outputs/k-next-20261020/convex/orientation_bonus_cover.log}
\end{quote}
The strongest default cover verifier is
\begin{quote}\small\ttfamily\raggedright
\nolinkurl{/Users/alingalatan/Documents/Codex/2026-10-10/plea/outputs/k-next-20261020/convex/check_orientation_bonus.cpp}
\end{quote}
Its completed receipt establishes Theorem~2, including the sector penalty,
together with the analytic reductions and exact integration inputs.  The
unstrengthened cover passed separately as well.  This is a theorem for the
stated Gaussian-factor class, not a complete row envelope or an improvement
of the global Grothendieck endpoint.

\subsection{A Gaussian-marginal obstruction to a law-only shared comparison}
\label{sec:gaussian-marginal-shared-obstruction}
This is not a counterexample to Conjecture K: the example below is not a pure cubic chaos. It shows why knowing the exact radial law, charge, first-chaos orthogonality and the covariance profile is still insufficient.

Let $L=Z(a+b|Z|^2)$, with $a>0$ and $b\ge0$, and suppose
\[
g(v)=2\mathbb E|L+\sqrt vG|-v/\epsilon
\]
has a maximizer $v_*>0$. Write $S=g(v_*)$. For a standard proper complex Gaussian $X$, put
\[
R=|X|^2,\qquad F_2(r)=1-(r+1)e^{-r},\qquad
Q=\sqrt{v_*}\,X\exp(2\pi iF_2(R)).
\]
Then $Q$ has exactly the proper Gaussian law of variance $v_*$. Indeed its modulus is unchanged and its uniform phase is translated by a function of the modulus. Moreover,
\[
\frac{\mathbb E[Q\overline X]}{\sqrt{v_*}}
 =\int_0^\infty r e^{-r}e^{2\pi iF_2(r)}\,dr
 =\int_0^1 e^{2\pi iy}\,dy=0.
\]
Charge one also gives $\mathbb EQ=\mathbb EQ^2=\mathbb E[QX]=0$. Consequently $Q+\alpha(T-1)X$, conditionally on the visible radius $T$, has variance $v_*+|\alpha|^2(T-1)^2$ and zero pseudovariance: exactly the covariance profile of the Gaussian A theorem.

Nevertheless, some arbitrarily small A perturbations increase the row above $S$. To see this, use the completely monotone scalar potential
\[
F_A(u)=\mathbb E[(T-1)|L+x|],\qquad |x|^2=u,
\]
and set $w(u)=-F_A'(u)$. Its positive Laplace mixture implies that $w$ is positive and strictly decreasing. Define
\[
J=\mathbb E[R e^{2\pi iF_2(R)}w(v_*R)].
\]
Differentiation of the honest two-sign payoff gives
\[
D_\alpha\Lambda\big|_{\alpha=0}
=-4\sqrt{v_*}\,\Re(\alpha\overline J).
\]
The price has no linear term because $\mathbb E[Q\overline X]=0$. A change of variables shows
\[
\Im J=\int_0^{1/2}\sin(2\pi y)
 \bigl[w(v_*F_2^{-1}(y))-w(v_*F_2^{-1}(1-y))\bigr]\,dy>0.
\]
Thus $\alpha=-it$, $t>0$, yields
\[
\Lambda(Q-it(T-1)X)=S+4\sqrt{v_*}(\Im J)t+o(t)>S
\]
for all sufficiently small positive $t$.

The positive-maximizer assumption holds at the peak profile. In fact
\[
g'(0)=\tfrac12\mathbb E|L|^{-1}-\epsilon^{-1},\qquad
\mathbb E|L|^{-1}\ge\frac{\sqrt\pi}{a+b/2}.
\]
At $a=699/1250$, $b=179/1250$, $\epsilon=1891/2500$, the rational estimate $\sqrt\pi>17/10$ already gives $g'(0)>0$. Continuity and $g(v)\to-\infty$ prove existence of a positive maximizer.

The precise missing assumption is spectral. With real standard Gaussian coordinates $X=(x+iy)/\sqrt2$, direct differentiation gives
\[
\mathbb E\|DQ\|^2
 =v_*\left(1+\frac{64\pi^2}{81}\right)>3v_*.
\]
Here $\phi'(r)=2\pi r e^{-r}$ and $\mathbb E[R^4e^{-2R}]=8/81$. Charge and first-chaos orthogonality allow the bidegrees $(2,1),(3,2),\ldots$; they do not force $(2,1)$. A pure cubic has Dirichlet energy exactly $3v_*$. In particular this example does not pass the actual chaos Stein inputs\stale{22 Oct}{proved in Lean for the Gaussian law: \code{gaussian\_steinInputs}} just because its marginal admits a constant abstract Stein kernel.

Every remote radial test, and every test depending only on adding independent Gaussian noise, sees exactly a Gaussian in this example. Such information therefore cannot control shared A content without a genuine joint fixed-chaos identity. No Grothendieck endpoint is improved by this obstruction.
